Podsumowując, niniejsza praca polegała na stworzeniu serwisu z ogłoszeniami o sprzedaży samochodów, który skupiałby się na poprawieniu obecnej sytuacji kupujących, na rynku wtórnym sprzedaży pojazdów mechanicznych.
Projekt został zrealizowany w oparciu o architekturę klient-serwer z komunikacją zgodnie z architekturą REST wykorzystującą protokół HTTP.
Został zaimplementowany model ról i uprawnień, który odpowiednio udostępnia funkcjonalności systemu użytkownikom, zależnie od statusu ich autentykacji oraz posiadanych ról.
Wykorzystano szereg różnorodnych technologii w ramach wytworzenia rozwiązania, które były zarówno nowoczesne, jak i sprawdzone w praktyce.

Serwer jest aplikacją obsługującą żądania zgodnie z modelem warstwowym, wykorzystując JSON jako główny format przetwarzania danych.
Najważniejszym elementem rozwiązania jest system przeliczania reputacji sprzedawcy, który ma wpływ na pozycjonowanie jego ofert w wynikach wyszukiwania.
Część kliencka to aplikacja internetowa z wieloma widokami pozwalającymi na korzystanie z możliwości systemu w intuicyjny sposób.
Dużą uwagę przywiązano do warstwy wizualnej, wygody korzystania z formularzy oraz wsparcia dla urządzeń mobilnych.

Do zrealizowanych wymagań można zaliczyć również zamieszczanie opinii pod ofertami, ogólnodostępne archiwum ogłoszeń sprzedających oraz system zgłaszania i moderowania treścią na platformie.
Nie zostało natomiast w pełni wdrożone pozycjonowanie ofert bazujące na warunkach takich jak, czy użytkownik był na oględzinach pojazdu.

Projekt posiada liczne ścieżki dalszego rozwoju.
Jedną z nich jest rozwój algorytmu pozycjonowania ofert, możliwie oparty o algorytmy sztucznej inteligencji, pod które dokumentacyjne bazy danych są zoptymalizowane, takie jak ta wykorzystana w projekcie.
Kolejnym jest dodanie brakującego wymagania z wbudowanym systemem umawiania się ze sprzedawcą na oględziny sprzedawanego samochodu.
Pomysłem wartym do zwalidowania jest również pierwotny pomysł nakładki, na przykład w postaci rozszerzenia dla przeglądarek lub platformy wykorzystującej roboty indeksujące zewnętrzne serwisy.

Wszystkie potencjalne kierunki rozwoju znacząco zwiększą funkcjonalności aplikacji, a co za tym idzie atrakcyjność dla przyszłych użytkowników.
Wraz z postępem aplikacji na rynku, należy ją odpowiednio skalować, co opierałoby się wstępnie na zwiększaniu samej bazy danych pod względem pojemności i prędkości.
Całość projektu jest zoptymalizowana pod daleko idący rozwój.
